\chapter*{GENERAL INTRODUCTION}
\markboth{\MakeUppercase{GENERAL INTRODUCTION}}{}
\addcontentsline{toc}{chapter}{GENERAL INTRODUCTION}
\adjustmtc
\thispagestyle{MyStyle}

In industrial manufacturing, a Request for Quotation (RFQ) is rarely a simple commercial exchange. It is the entry point of a long, multi-departmental process that mobilises engineering, production, quality, and procurement expertise over weeks or months, and whose outcome can determine whether a project is won, lost, or accepted under unfavourable terms. The work presented in this report addresses not the production of the quotation itself, but the structural fragility of the process that surrounds it --- the coordination, the visibility, the historical learning, and the early surfacing of feasibility risk that together separate a disciplined estimation organisation from one that relies on individual diligence to compensate for the absence of system-level support.

This end-of-studies project was conducted as an internship at \hostCompany, a Tunisia-based consulting firm specialising in cloud architecture, microservices, artificial intelligence, and data engineering. The practical anchor of the work was a five-day on-site audit at \clientCompany, a major industrial conglomerate based in Dammam, Saudi Arabia, and its subsidiary GHI (Gulf Heavy Industries Co.) --- one of the largest specialised heavy steel fabrication companies in the Middle East. The audit produced twelve documented pain points clustered into five structural families: coordination, visibility, historical learning, feasibility and supplier risk, and tooling fragility.

The client's initial mandate had been the construction of an artificial-intelligence system capable of automating the production of the Bill of Quantities. The audit findings pointed in a different direction: the senior estimation expertise embedded in the workbook was not extractable as training data, and the asymmetric cost of error in industrial quotation left no tolerable failure mode for an automated pricing system. Rather than weaken the original goal into a partial assistant whose value would be marginal and whose risk would not, the project was reframed: build a platform that strengthens every part of the process that does not depend on irreducible human expertise, and respect the part that does.

The platform that resulted is named \textit{Itqan}, after the Arabic word for mastery, precision, and the act of doing work properly and to completion. Concretely, Itqan is a web platform that tracks RFQs across their lifecycle, centralises stages, artifacts and reminders, provides portfolio dashboards, parses RFQ/workbook documents deterministically, and exposes a controlled AI copilot that answers only from authorised platform evidence. It is organised around a four-pillar target architecture --- Workflow and Tracking, Communication Automation, Executive Intelligence and Business Intelligence, and Historical Learning and Prediction --- and is augmented by a cross-cutting conversational layer, the \textit{RFQ Copilot}, whose central engineering contribution is the trust-boundary architecture that disciplines its behaviour under the cost-of-error conditions of the estimation domain. The first pillar is substantially delivered; the second and third are partial; the fourth is delivered at the foundation level through deterministic document parsing, with predictive components left as future scope.

The report develops the project in six chapters. Chapter~1 establishes the host organisation, the industrial context, the audit findings, and the strategic pivot. Chapter~2 positions the platform against the existing technological and academic landscape. Chapter~3 formalises the actors, requirements, use cases, and methodology. Chapter~4 develops the conceptual and technical architecture, with the trust-boundary architecture of the copilot as its central innovation contribution. Chapter~5 presents the implementation across the delivered platform components: the manager service, the copilot service, the intelligence service, and the user interface. Chapter~6 presents the validation evidence verifying that the architectural invariants hold under measured behaviour. The general conclusion revisits the contributions and the direction of continued evolution.
